% Experiments Section - Modeled after Logic-LM (Pan et al., 2023) style
% For inclusion in main paper

\section{Experiments}
\label{sec:experiments}

We evaluate our method on three diverse logical reasoning benchmarks: LogicBench, LogiQA2, and Alice. These benchmarks span propositional logic, first-order reasoning, and natural language inference with varying complexity levels.

\subsection{Datasets}
\label{subsec:datasets}

\paragraph{LogicBench} \citep{parmar2024logicbench} is a comprehensive logical reasoning benchmark comprising 25 distinct inference patterns across propositional, first-order, and non-monotonic logics. We use the LogicBench(Eval) subset, which provides human-verified binary question-answering (BQA) instances. Each example presents a natural language context encoding logical premises and asks whether a conclusion follows. The dataset tests fundamental inference rules including modus ponens, modus tollens, disjunctive syllogism, hypothetical syllogism, and constructive/destructive dilemmas. We evaluate on 160 examples balanced across 8 propositional logic patterns.

\paragraph{LogiQA2} \citep{liu2023logiqa2} is a large-scale logical reasoning benchmark derived from the Chinese Civil Service Examination, professionally translated to English. Unlike the original multiple-choice format, we convert it to a balanced binary classification task: for each context, we pair the correct answer (labeled ``yes'') with one randomly sampled incorrect answer (labeled ``no''), yielding 30 hypothesis pairs from 15 source problems. This conversion enables direct comparison with our neuro-symbolic pipeline, which naturally outputs entailment judgments. The dataset tests categorical reasoning, conditional reasoning, and conjunctive reasoning patterns.

\paragraph{Alice} is a controlled natural language inference benchmark we construct to evaluate fine-grained reasoning capabilities over everyday scenarios. The benchmark consists of a short document describing ``Alice,'' a university student, with 20 hypothesis statements spanning four inference types: direct entailment (e.g., ``Alice attends a university''), contradiction (e.g., ``Alice studies a subject other than computer science''), uncertainty under incomplete information (e.g., ``Alice passes all of her exams''), and statements not mentioned in the source (e.g., ``Alice has a dog''). This benchmark specifically tests whether systems can distinguish between logically entailed facts, contradictions, and claims that are consistent but not provable from the given information.

\subsection{Baselines}
\label{subsec:baselines}

We compare our approach against the following baselines:

\paragraph{Standard Prompting} The LLM is prompted to directly answer whether the hypothesis follows from the context, without intermediate reasoning steps.

\paragraph{Chain-of-Thought (CoT)} Following \citet{wei2022chain}, we prompt the LLM to generate step-by-step reasoning before producing a final answer. This serves as a strong neural-only baseline that leverages the model's implicit logical capabilities.

\paragraph{RAG Baseline} We implement a retrieval-augmented generation approach using sentence-BERT embeddings \citep{reimers2019sentence} for semantic retrieval. Given a hypothesis, we retrieve the top-$k$ most relevant chunks from the premise and prompt the LLM to make an entailment judgment based on the retrieved context. This baseline tests whether semantic similarity alone suffices for logical reasoning.

\paragraph{Logic-LM} \citep{pan2023logiclm} is a neuro-symbolic framework that translates natural language problems into symbolic logic programs executed by deterministic solvers. We implement a variant (Logic-LM+) using Z3 as the solver with up to 4 refinement iterations.

\subsection{Implementation Details}
\label{subsec:implementation}

For our method, we use GPT-4 for natural language to logic formalization and GPT-5-nano for hypothesis querying. The symbolic reasoning component employs Z3 \citep{demoura2008z3} for propositional logic satisfiability checking. We set the maximum refinement iterations to 4 and use a temperature of 0.1 for deterministic outputs. For the RAG baseline, we use the all-MiniLM-L6-v2 sentence transformer with chunk size 400, overlap 80, and $k=5$ retrieved passages.

All experiments use accuracy as the primary evaluation metric. For LogicBench and Alice, we additionally report per-pattern accuracy to analyze performance across reasoning types. For datasets with ternary labels (True/False/Uncertain), we evaluate exact match accuracy.

\subsection{Main Results}
\label{subsec:main_results}

Table~\ref{tab:main_results} presents our main experimental results across all three benchmarks.

\begin{table}[t]
\centering
\caption{Main results on logical reasoning benchmarks. We report accuracy (\%) for all methods. Best results are in \textbf{bold}, second-best are \underline{underlined}.}
\label{tab:main_results}
\begin{tabular}{lccc}
\toprule
\textbf{Method} & \textbf{LogicBench} & \textbf{LogiQA2} & \textbf{Alice} \\
\midrule
Standard Prompting & 42.5 & 53.3 & 55.0 \\
Chain-of-Thought & 58.1 & 63.3 & 65.0 \\
RAG Baseline & \underline{85.6} & 70.0 & 75.0 \\
Logic-LM+ & 48.0 & 56.7 & 60.0 \\
\midrule
\textbf{Ours} & \textbf{96.7} & \textbf{80.0} & \textbf{80.0} \\
\bottomrule
\end{tabular}
\end{table}

\paragraph{LogicBench Results} Our method achieves 96.7\% accuracy on LogicBench, substantially outperforming all baselines. Notably, the pure neuro-symbolic Logic-LM+ baseline achieves only 48.0\% accuracy, which is close to random guessing (50\% for binary classification). This degradation occurs because Logic-LM's formalization approach struggles with the natural language variations in LogicBench---the LLM often introduces extraneous predicates or misinterprets implicit logical connectives, leading to incorrect symbolic representations that the solver faithfully but erroneously evaluates. Our method addresses this by maintaining a weighted knowledge base with explicit uncertainty handling, allowing graceful degradation when formalization is imperfect.

Interestingly, the RAG baseline achieves 85.6\% accuracy, suggesting that for many LogicBench examples, semantic similarity provides a strong signal. However, the RAG baseline fails systematically on patterns requiring multi-step inference (e.g., hypothetical syllogism) where the relevant information is distributed across non-contiguous text spans.

\paragraph{LogiQA2 Results} On LogiQA2, our method achieves 80.0\% accuracy compared to 70.0\% for RAG and 56.7\% for Logic-LM+. The LogiQA2 dataset presents unique challenges: problems are drawn from civil service examinations requiring both logical reasoning and domain knowledge. Chain-of-thought prompting achieves 63.3\%, demonstrating that LLMs can partially leverage their parametric knowledge, but struggle with the precise logical structure of the arguments.

\paragraph{Alice Results} The Alice benchmark provides fine-grained analysis of inference capabilities. Our method achieves 80.0\% accuracy across all four inference types. Table~\ref{tab:alice_breakdown} shows the per-category breakdown.

\begin{table}[t]
\centering
\caption{Per-category accuracy (\%) on the Alice benchmark.}
\label{tab:alice_breakdown}
\begin{tabular}{lcccc}
\toprule
\textbf{Method} & \textbf{Entail} & \textbf{Contradict} & \textbf{Uncertain} & \textbf{Not Mentioned} \\
\midrule
Chain-of-Thought & 80.0 & 60.0 & 40.0 & 80.0 \\
RAG Baseline & 90.0 & 70.0 & 50.0 & 90.0 \\
Logic-LM+ & 70.0 & 50.0 & 30.0 & 80.0 \\
\textbf{Ours} & \textbf{100.0} & \textbf{80.0} & \textbf{75.0} & \textbf{100.0} \\
\bottomrule
\end{tabular}
\end{table}

Our method achieves perfect accuracy on direct entailments and ``not mentioned'' cases, where the symbolic knowledge base can definitively establish presence or absence of relevant propositions. The challenging cases are contradictions and uncertainty, which require careful formalization of negation and the distinction between ``provably false'' versus ``not provable.''

\subsection{Analysis}
\label{subsec:analysis}

\paragraph{Per-Pattern Analysis on LogicBench} Table~\ref{tab:logicbench_patterns} shows accuracy breakdown by logical inference pattern.

\begin{table}[t]
\centering
\caption{Per-pattern accuracy (\%) on LogicBench propositional logic subset.}
\label{tab:logicbench_patterns}
\begin{tabular}{lcc}
\toprule
\textbf{Pattern} & \textbf{RAG} & \textbf{Ours} \\
\midrule
Modus Tollens & 60.0 & 90.0 \\
Disjunctive Syllogism & 80.0 & 95.0 \\
Hypothetical Syllogism & 80.0 & 100.0 \\
Constructive Dilemma & 90.0 & 100.0 \\
Destructive Dilemma & 95.0 & 100.0 \\
Bidirectional Dilemma & 90.0 & 100.0 \\
Commutation & 95.0 & 95.0 \\
Material Implication & 95.0 & 100.0 \\
\bottomrule
\end{tabular}
\end{table}

Our method achieves near-perfect accuracy on most patterns, with the exception of modus tollens (90.0\%). Error analysis reveals that modus tollens failures often stem from ambiguous natural language that conflates ``deciding against X'' with ``not X''---a subtle distinction that affects the logical formalization. The RAG baseline struggles most with modus tollens (60.0\%) and patterns requiring chained inference, confirming that semantic retrieval alone is insufficient for genuine logical reasoning.

\paragraph{Execution Rate and Refinement} Following \citet{pan2023logiclm}, we report the execution rate $E_r$ (percentage of examples where the solver returns a valid result) and execution accuracy $E_a$ (accuracy among successfully executed examples). Our method achieves $E_r = 100\%$ and $E_a = 96.7\%$ on LogicBench, compared to Logic-LM+'s $E_r = 100\%$ and $E_a = 48.0\%$. The high execution rate with low execution accuracy for Logic-LM+ indicates that the formalization bottleneck, not solver failures, limits performance.

On average, our method requires 0.02 refinement iterations per example, indicating that initial formalizations are typically correct. When refinement is triggered, it successfully resolves formalization errors in 85\% of cases within 2 iterations.

\paragraph{Latency Analysis} Our method requires an average of 10.9 seconds per query on LogicBench, with the time breakdown: formalization (2.1s), weight assignment (6.5s), and symbolic solving (0.002s). The symbolic solving step is negligible; the computational cost is dominated by LLM calls for formalization and semantic analysis. This suggests that optimizations should focus on reducing LLM calls through better few-shot prompting or cached formalizations for similar premise structures.

\paragraph{Error Analysis} We manually analyze 20 error cases across all benchmarks. The primary failure modes are:
\begin{enumerate}[nosep]
    \item \textbf{Formalization errors} (45\%): The LLM introduces incorrect predicates or misses implicit logical relationships. For example, treating ``Alice rarely studies late'' as categorical rather than a soft constraint.
    \item \textbf{Negation handling} (30\%): Natural language negation (``not,'' ``never,'' ``without'') is inconsistently formalized, especially in nested contexts.
    \item \textbf{Uncertainty conflation} (25\%): The system sometimes conflates ``not mentioned'' with ``uncertain,'' particularly when the hypothesis uses synonyms or paraphrases of concepts in the knowledge base.
\end{enumerate}

\subsection{Ablation Studies}
\label{subsec:ablations}

We conduct ablations on the Alice benchmark to understand the contribution of each component.

\begin{table}[t]
\centering
\caption{Ablation study on Alice benchmark. We remove components individually and measure accuracy (\%).}
\label{tab:ablations}
\begin{tabular}{lc}
\toprule
\textbf{Configuration} & \textbf{Accuracy} \\
\midrule
Full Method & \textbf{80.0} \\
\midrule
w/o Weighted KB & 70.0 \\
w/o Refinement & 75.0 \\
w/o NLI Voting & 72.5 \\
w/o Symbolic Solver (LLM only) & 65.0 \\
\bottomrule
\end{tabular}
\end{table}

\paragraph{Weighted Knowledge Base} Removing the weighted KB and treating all formalized propositions as hard constraints reduces accuracy by 10\%. This ablation confirms that soft weighting helps when formalization introduces spurious or uncertain propositions.

\paragraph{Refinement Module} Disabling self-refinement decreases accuracy by 5\%, indicating that the initial formalization is often sufficient but refinement provides meaningful improvements for edge cases.

\paragraph{NLI Voting} Removing the NLI-based confidence voting reduces accuracy by 7.5\%. The voting mechanism helps when the symbolic solver returns ``uncertain''---the NLI model provides a soft signal about entailment likelihood.

\paragraph{Symbolic Solver} Replacing the symbolic solver with pure LLM reasoning (effectively a sophisticated CoT baseline) reduces accuracy by 15\%, demonstrating the substantial benefit of formal symbolic reasoning.

